\documentclass[10pt]{article}

\usepackage[
  paperwidth=16cm,
  paperheight=9cm,
  margin=0cm
]{geometry}
\usepackage[table]{xcolor}
\usepackage{array}
\usepackage{microtype}

\pagestyle{empty}
\setlength{\parindent}{0pt}
\setlength{\fboxrule}{0.35pt}

\definecolor{Carbon}{HTML}{0A0A0A}
\definecolor{Cream}{HTML}{FCFCFC}
\definecolor{MutedGold}{HTML}{A16207}
\definecolor{SlateGrey}{HTML}{444444}
\definecolor{DeepGrey}{HTML}{151515}
\definecolor{SoftGrey}{HTML}{B8B8B8}

\newcounter{slideno}
\newcommand{\repoid}{github.com/SheikhMohsin9311/Benchmarking-and-Performance}
\newcommand{\presenter}{Sheikh Mohsin}
\newcommand{\slidefooter}[1]{%
  \vfill
  {\color{MutedGold}\small\hfill #1\,/\,\pageref{LastSlide}}%
}

\newcommand{\titleslide}{%
  \stepcounter{slideno}
  \pagecolor{Carbon}
  \color{Cream}
  \begin{minipage}[t][8.72cm][c]{16cm}
    \hspace*{1.05cm}
    \begin{minipage}{13.9cm}
      {\color{MutedGold}\Large Paper Presentation}

      \vspace{0.45cm}
      {\Huge\bfseries How to Build a Benchmark}

      \vspace{0.35cm}
      {\large\itshape Kistowski, Arnold, Lange, Henning, Huppler, and Cao}

      \vspace{0.65cm}
      {\color{SlateGrey}\large ICPE 2015 \quad | \quad Benchmarking and Performance}

      \vspace{0.9cm}
      {\color{SlateGrey}\rule{5.4cm}{0.4pt}}

      \vspace{0.55cm}
      {\color{SlateGrey}A practitioner's vocabulary for benchmark quality}

      \vspace{0.75cm}
      {\color{MutedGold}\small \presenter}

      \vspace{0.12cm}
      {\color{SlateGrey}\small \repoid}
    \end{minipage}
  \end{minipage}
  \newpage
}

\newcommand{\sectionslide}[2]{%
  \stepcounter{slideno}
  \begin{minipage}[t][8.72cm][c]{16cm}
    \hspace*{1.05cm}
    \begin{minipage}{13.9cm}
      {\color{MutedGold}\Large #1}

      \vspace{0.45cm}
      {\Huge\bfseries #2}

      \vspace{0.65cm}
      {\color{SlateGrey}\rule{4.2cm}{0.4pt}}
    \end{minipage}
  \end{minipage}
  \newpage
}

\newcommand{\contentslide}[2]{%
  \stepcounter{slideno}
  \begin{minipage}[t][8.72cm][t]{16cm}
    \vspace*{0.72cm}

    \hspace*{0.82cm}
    \begin{minipage}{14.35cm}
      {\color{Cream}\Large\bfseries #1}

      \vspace{0.18cm}
      {\color{SlateGrey}\rule{3.2cm}{0.35pt}}

      \vspace{0.48cm}
      #2
      \slidefooter{\arabic{slideno}}
    \end{minipage}
  \end{minipage}
  \newpage
}

\newcommand{\tagline}[1]{{\color{MutedGold}\bfseries #1}}
\newcommand{\muted}[1]{{\color{SlateGrey}#1}}
\newcommand{\soft}[1]{{\color{SoftGrey}#1}}
\newcommand{\emphasis}[1]{{\color{Cream}\bfseries #1}}
\newcommand{\callout}[1]{%
  \begingroup
  \setlength{\fboxsep}{0.2cm}%
  \fcolorbox{MutedGold}{Carbon}{%
    \begin{minipage}{\dimexpr\linewidth-2\fboxsep-2\fboxrule\relax}
      \vspace{0.12cm}
      {\color{Cream}\itshape #1}
      \vspace{0.12cm}
    \end{minipage}%
  }
  \endgroup
}
\newcommand{\darkbox}[1]{%
  \begingroup
  \setlength{\fboxsep}{0.2cm}%
  \fcolorbox{SlateGrey}{DeepGrey}{%
    \begin{minipage}{\dimexpr\linewidth-2\fboxsep-2\fboxrule\relax}
      \vspace{0.1cm}
      #1
      \vspace{0.1cm}
    \end{minipage}%
  }
  \endgroup
}
\newcommand{\twocol}[2]{%
  \begin{minipage}[t]{6.75cm}#1\end{minipage}
  \hfill
  \begin{minipage}[t]{6.75cm}#2\end{minipage}
}
\newcommand{\threecol}[3]{%
  \begin{minipage}[t]{4.35cm}#1\end{minipage}
  \hfill
  \begin{minipage}[t]{4.35cm}#2\end{minipage}
  \hfill
  \begin{minipage}[t]{4.35cm}#3\end{minipage}
}
\newcommand{\metric}[2]{%
  \centering{\color{MutedGold}\Huge\bfseries #1}\\[0.25em]
  {\color{SlateGrey}#2}
}

\begin{document}

\pagecolor{Carbon}
\color{Cream}

\titleslide

\contentslide{One-Sentence Reading}{
  \vfill
  \callout{An insider account from SPEC and TPC committee members on how standardized benchmarks are defined, governed, validated, and made usable.}
  \vspace{0.65cm}

  \twocol{
    \tagline{What the paper gives us}
    \begin{itemize}
      \item A taxonomy of benchmark forms
      \item A vocabulary for quality
      \item A map of design trade-offs
    \end{itemize}
  }{
    \tagline{Why it matters}
    \begin{itemize}
      \item Benchmarks turn measurements into claims
      \item Claims need trust, context, and rules
      \item The process is part of the result
    \end{itemize}
  }
}

\contentslide{Why This Paper Exists}{
  Benchmark consortia such as \emphasis{SPEC} and \emphasis{TPC} accumulate practical knowledge under confidentiality agreements.

  \vspace{0.35cm}
  This paper opens that process enough for researchers to see how serious standardized benchmarks are shaped.

  \vspace{0.5cm}
  \threecol{
    \tagline{SPEC}

    \soft{SPEC CPU, SPECpower, SERT}
  }{
    \tagline{TPC}

    \soft{TPC-C, TPC-H, TPC-DS}
  }{
    \tagline{Audience}

    \soft{Benchmark designers, users, and critics}
  }
}

\sectionslide{Part 1}{Definitions}

\contentslide{Benchmark or Rating Tool?}{
  \twocol{
    \darkbox{
      \tagline{Benchmark}

      \vspace{0.14cm}
      A standard tool for \emphasis{competitive} comparison of systems or components.
    }

    \vspace{0.28cm}
    \begin{itemize}
      \item Product comparison
      \item Published ranking
      \item Strict comparability pressure
    \end{itemize}
  }{
    \darkbox{
      \tagline{Rating tool}

      \vspace{0.14cm}
      A standard tool for \emphasis{non-competitive} evaluation, regulation, research, or system improvement.
    }

    \vspace{0.28cm}
    \begin{itemize}
      \item Diagnosis over ranking
      \item Broader experimental use
      \item Example: SPEC SERT
    \end{itemize}
  }
}

\contentslide{Three Benchmark Forms}{
  \threecol{
    \tagline{Specification-based}

    \vspace{0.16cm}
    Defines the problem, inputs, and expected outcomes. The submitter builds the implementation.

    \vspace{0.2cm}
    \soft{Flexible, but compliance is harder.}
  }{
    \tagline{Kit-based}

    \vspace{0.16cm}
    Provides the implementation as part of the official benchmark artifact.

    \vspace{0.2cm}
    \soft{Easy to run, but less open to innovation.}
  }{
    \tagline{Hybrid}

    \vspace{0.16cm}
    Provides most of the benchmark while allowing selected parts to vary.

    \vspace{0.2cm}
    \soft{A compromise between control and flexibility.}
  }
}

\contentslide{The Design Thesis}{
  \vfill
  \callout{No single workload can be strong across every quality dimension. Benchmark design is a controlled compromise.}

  \vspace{0.65cm}
  \twocol{
    \tagline{Narrow and faithful}
    \begin{itemize}
      \item High relevance in one domain
      \item Lower breadth of applicability
      \item Easier to interpret locally
    \end{itemize}
  }{
    \tagline{Broad and portable}
    \begin{itemize}
      \item More systems can participate
      \item Less precise for any one scenario
      \item Easier to standardize widely
    \end{itemize}
  }
}

\sectionslide{Part 2}{Quality Criteria}

\contentslide{The Five Criteria}{
  \vspace{0.1cm}
  \threecol{
    \metric{1}{Relevance}
  }{
    \metric{2}{Reproducibility}
  }{
    \metric{3}{Fairness}
  }

  \vspace{0.75cm}
  \twocol{
    \metric{4}{Verifiability}
  }{
    \metric{5}{Usability}
  }
}

\contentslide{1. Relevance}{
  \callout{Does the benchmark behavior correlate with what consumers of the result actually care about?}

  \vspace{0.38cm}
  \twocol{
    \tagline{Two dimensions}
    \begin{itemize}
      \item Breadth of applicability
      \item Degree of domain fidelity
      \item Scalability for server workloads
    \end{itemize}
  }{
    \tagline{Design move}
    \begin{itemize}
      \item Start from intended use
      \item Identify result consumers
      \item Work backward from decisions they will make
    \end{itemize}
  }
}

\contentslide{2. Reproducibility}{
  \callout{Reproducibility has two faces: repeated runs on the same setup, and equivalent runs by an independent tester.}

  \vspace{0.35cm}
  \twocol{
    \tagline{Variability sources}
    \begin{itemize}
      \item Thread scheduling and JIT effects
      \item Disk layout and network contention
      \item Power management and temperature
    \end{itemize}
  }{
    \tagline{Control mechanisms}
    \begin{itemize}
      \item Long enough steady-state runs
      \item Multiple run requirements
      \item Full hardware, software, and tuning disclosure
    \end{itemize}
  }
}

\contentslide{3. Fairness}{
  \callout{Fairness means systems compete on their merits, not on benchmark loopholes.}

  \vspace{0.35cm}
  \twocol{
    \tagline{Too restrictive}
    \begin{itemize}
      \item Excludes legitimate setups
      \item Reduces participation
      \item Narrows result usefulness
    \end{itemize}
  }{
    \tagline{Too permissive}
    \begin{itemize}
      \item Rewards unrealistic configurations
      \item Pollutes result pools
      \item Encourages cherry-picked comparisons
    \end{itemize}
  }
}

\contentslide{4. Verifiability}{
  \callout{Industry benchmark results are often produced by interested vendors; trust has to be engineered into the benchmark.}

  \vspace{0.35cm}
  \begin{itemize}
    \item Runtime self-validation checks that the workload is behaving legally.
    \item Functional output checks catch invalid optimizations.
    \item Benchmark-collected configuration data is stronger than user-declared data.
    \item Redundant result metadata can reveal inconsistencies in reported scores.
  \end{itemize}
}

\contentslide{5. Usability}{
  \callout{A benchmark that is too costly, complex, or fragile to run will not become a useful standard.}

  \vspace{0.35cm}
  \twocol{
    \tagline{Cost barrier}
    \begin{itemize}
      \item Large benchmarks may require specialized systems
      \item TPC-C examples can reach enterprise-scale cost
      \item High cost reduces participation
    \end{itemize}
  }{
    \tagline{Usability mechanisms}
    \begin{itemize}
      \item Self-validation
      \item Load-and-go benchmark kits
      \item Automated configuration collection
    \end{itemize}
  }
}

\sectionslide{Part 3}{Tensions}

\contentslide{Criteria Pull Against Each Other}{
  \twocol{
    \tagline{Relevance vs. reproducibility}

    \soft{Realistic workloads often introduce more variability.}

    \vspace{0.35cm}
    \tagline{Relevance vs. usability}

    \soft{Representative workloads can demand expensive systems.}

    \vspace{0.35cm}
    \tagline{Fairness vs. relevance}

    \soft{Run rules can exclude legitimate real-world configurations.}
  }{
    \tagline{Verifiability vs. usability}

    \soft{More validation and reporting adds friction.}

    \vspace{0.35cm}
    \tagline{Specificity vs. breadth}

    \soft{High domain fidelity usually narrows applicability.}

    \vspace{0.35cm}
    \callout{The point is not to eliminate trade-offs. The point is to make them deliberate.}
  }
}

\contentslide{The Governance Point}{
  \vfill
  \callout{Consensus is not just bureaucracy. In standardized benchmarking, governance is one of the mechanisms that makes fairness credible.}

  \vspace{0.6cm}
  \begin{itemize}
    \item Committees force trade-offs to be negotiated across stakeholders.
    \item Run rules define what counts as a valid result.
    \item Fair-use policies govern how results may be compared publicly.
  \end{itemize}
}

\sectionslide{Part 4}{For Our Work}

\contentslide{How This Maps to Our Project}{
  \twocol{
    \tagline{Early decision}
    \begin{itemize}
      \item Are we building a benchmark or a rating tool?
      \item Are we ranking systems or diagnosing behavior?
      \item Who consumes the result?
    \end{itemize}
  }{
    \tagline{SUT definition}
    \begin{itemize}
      \item Hardware and firmware
      \item OS, compiler, drivers, libraries
      \item Runtime settings and tuning knobs
    \end{itemize}
  }
}

\contentslide{Modern Pressure Points}{
  \begin{itemize}
    \item \tagline{ML nondeterminism:} random seeds, data shuffling, and floating-point non-associativity complicate reproducibility.
    \item \tagline{Cloud systems:} the boundary of the system under test becomes harder to describe.
    \item \tagline{Energy metrics:} power management and temperature make measurement variability worse.
    \item \tagline{MLPerf-style divisions:} closed/open rules are fairness mechanisms, but they also shape interpretation.
  \end{itemize}
}

\contentslide{Final Takeaways}{
  \vspace{0.2cm}
  \begin{enumerate}
    \item Define the intended use before choosing workloads.
    \item Treat workload, metric, run rules, validation, and reporting as one system.
    \item Make trade-offs explicit; they are not defects if they are justified.
    \item Prefer reproducible evidence over isolated headline numbers.
  \end{enumerate}

  \vspace{0.5cm}
  \callout{A strong benchmark does not merely measure performance. It creates evidence other people can interpret, challenge, and reproduce.}
}

\contentslide{Discussion}{
  \Large
  \begin{enumerate}
    \item Which benchmark type fits our project: specification, kit, or hybrid?
    \item Which quality criterion is most at risk for us?
    \item What should be standardized first: workload, metric, environment, or reporting?
  \end{enumerate}

  \vspace{0.55cm}
  {\normalsize\color{MutedGold}\presenter}

  \vspace{0.1cm}
  {\small\color{SlateGrey}\repoid}

  \label{LastSlide}
}

\end{document}
